\documentclass[12pt]{article}
\usepackage[utf8]{inputenc}
\usepackage[T1]{fontenc}
\usepackage[margin=1in]{geometry}
\usepackage{setspace}
\usepackage{xcolor}
\usepackage[colorlinks=true, urlcolor=blue!60!black, linkcolor=blue!60!black]{hyperref}
\setlength{\parindent}{2em}
\setlength{\parskip}{1.25em}
\renewcommand{\baselinestretch}{1.0}
\begin{document}
\begin{center}
{\LARGE \textit{United States v. Butler}, 297 U.S. 1 (1936)\par}
\vspace{0.8em}
{\large Case Brief --- Four-Page Edition\par}
\end{center}
\vspace{1em}

\section{Memory Jogger}

Striking down the Agricultural Adjustment Act's processing tax 6--3, the Court adopted Hamilton's view that the Spending Clause confers an independent power to tax and spend for the general welfare, yet held that Congress may not use that power to coerce or purchase compliance with regulation of local agricultural production reserved to the states (\href{https://supreme.justia.com/cases/federal/us/297/1/}{297 U.S. 1}).

\section{Detailed Case Facts}

During the Depression, collapsing farm prices left farmers' purchasing power far below that of the 1909--1914 base period. Title I of the Agricultural Adjustment Act of 1933 authorized the Secretary of Agriculture to contract with farmers to reduce acreage or production in exchange for rental or benefit payments. To fund the payments, the Act levied a processing tax on the first domestic processing of covered commodities, including cotton, set at the gap between current farm prices and a statutory fair exchange value, plus a floor tax on existing inventories, with proceeds earmarked for the program. In July 1933 the Secretary proclaimed cotton payments and fixed the tax rates.

\section{Procedural History}

The United States presented a tax claim to William M.\ Butler and the other receivers of Hoosac Mills, an insolvent Massachusetts cotton mill, who recommended disallowing it as unconstitutional. The district court ordered the taxes paid; the First Circuit reversed. The Supreme Court granted certiorari and, on January 6, 1936, affirmed.

\section{Judicial Votes}

Affirmed, \textbf{6--3}.
\begin{itemize}
\item \textbf{Majority:} Justice Roberts, joined by Chief Justice Hughes and Justices Van Devanter, McReynolds, Sutherland, and Butler.
\item \textbf{Concurring:} None.
\item \textbf{Dissenting:} Justice Stone, joined by Justices Brandeis and Cardozo.
\end{itemize}

\section{Holding}

(1) Unlike a general taxpayer under \textit{Massachusetts v.\ Mellon}, a processor resisting an earmarked exaction that is an integral part of a regulatory scheme has standing to challenge the whole plan. (2) Adopting the Hamiltonian and Story view, the Court held that Article I, \S~8, cl.~1 confers a substantive power to tax and spend for the general welfare, not confined to the other enumerated powers. (3) The AAA nonetheless violated the Tenth Amendment as a scheme to coerce and purchase compliance with federal regulation of local agricultural production.

\section{Analysis of Opinions}

\textbf{Majority (Roberts).} The exaction was ``an indispensable part in the plan of regulation,'' not a true tax for the support of government, so the receivers could challenge it. Roberts described judicial review as laying the Constitution beside the statute. With the Government not invoking the Commerce Clause, agriculture then being viewed as local, the Act rested on the taxing and spending power. Rejecting Madison's view that ``general welfare'' refers only to the enumerated powers, which would make the clause ``mere tautology,'' he sided with Hamilton and Story. But he declined to decide whether farm aid served the general welfare, holding instead, under the Tenth Amendment, \textit{McCulloch}'s pretext warning, and the \textit{Child Labor Tax Case}, that the Act regulated a reserved subject. The program was not truly voluntary, since nonparticipants faced ruin from subsidized competitors; and even voluntary compliance could not be bought, because Congress cannot purchase what it cannot compel. Otherwise Congress could subsidize its way into controlling wages, hours, and industry nationwide.

\textbf{Stone (dissenting).} Courts judge power, not wisdom. The processing tax itself compelled nothing; any influence came from the spending. Coercion was ``groundless speculation,'' because ``Threat of loss, not hope of gain, is the essence of economic coercion,'' and many cotton farmers declined to participate. Most tellingly, a broad spending power is meaningless if Congress cannot attach conditions to make the money serve its purpose, which would mean seeds given without a duty to plant them. The real checks are a national purpose, no coercion of reserved action, and the conscience of Congress and the Executive.

\textbf{Later developments.} In 1937, \textit{Steward Machine} and \textit{Helvering v.\ Davis} kept \textit{Butler}'s Hamiltonian reading while adopting Stone's deference. \textit{South Dakota v.\ Dole} (1987) synthesized five limits on conditional spending, and \textit{NFIB v.\ Sebelius} (2012) revived the coercion limit for the Medicaid expansion.

\section{Examples: Future Applications}

\begin{enumerate}
\item \textbf{Same-Side (Invalid): A Runoff Cap Backed by Medicaid.} States that refuse a nitrogen-runoff cap lose all existing Medicaid and highway funds, 18\% of their budgets. The condition is unrelated under \textit{Dole} and coercive under \textit{NFIB}.
\item \textbf{Same-Side (Standing): An Earmarked Dairy Fee.} A fee on processors funds only payments to farmers who cut output. Under \textit{Butler}, the processor may challenge the whole scheme despite \textit{Mellon}.
\item \textbf{Opposite-Side (Valid): A New, Optional Grant.} A new \$2 billion grant for runoff monitoring offers only ``hope of gain,'' with clear, germane conditions.
\item \textbf{Opposite-Side (Valid): Apprenticeship Stipends.} Direct payments from the general Treasury for job training fall within the Hamiltonian spending power and \textit{Helvering}'s deference.
\item \textbf{Fence-Sitter: Withholding 25\% of Highway Funds.} The penalty is larger than \textit{Dole}'s 5\% but far smaller than \textit{NFIB}'s, and \textit{NFIB} did not fix where pressure becomes compulsion.
\end{enumerate}

\section{Critique}

\textbf{Slot-machine jurisprudence and half-hearted Hamiltonianism.} Scholars have long mocked Roberts's claim that review merely lays the Constitution beside a statute, since neither the General Welfare Clause nor the Tenth Amendment defines coercion or general welfare; formalist rhetoric masked contested value judgments. And as Stone showed, adopting Hamilton's view makes little sense if Congress cannot attach conditions; Roberts never explained why an ongoing funding condition differs from a contract.

\textbf{Earmarking formalism.} Standing and the regulatory character of a tax turned on whether Congress earmarked revenue rather than routing it through the Treasury. Because money is fungible, critics call the distinction a drafting trap rather than a federalism safeguard.

\section{Key Quotations}

\begin{quote}
``The power of courts to declare a statute unconstitutional is subject to two guiding principles of decision which ought never to be absent from judicial consciousness. One is that courts are concerned only with the power to enact statutes, not with their wisdom. The other is that, while unconstitutional exercise of power by the executive and legislative branches of the government is subject to judicial restraint, the only check upon our own exercise of power is our own sense of self-restraint. For the removal of unwise laws from the statute books appeal lies not to the courts, but to the ballot and to the processes of democratic government.''
\end{quote}
\noindent --- Justice Stone, dissenting, 297 U.S.\ at 78--79.

\begin{quote}
``While, therefore, the power to tax is not unlimited, its confines are set in the clause which confers it, and not in those of \S\ 8 which bestow and define the legislative powers of the Congress. It results that the power of Congress to authorize expenditure of public moneys for public purposes is not limited by the direct grants of legislative power found in the Constitution.''
\end{quote}
\noindent --- Justice Roberts, Opinion of the Court, 297 U.S.\ at 66.

\begin{quote}
``When an act of Congress is appropriately challenged in the courts as not conforming to the constitutional mandate, the judicial branch of the Government has only one duty---to lay the article of the Constitution which is invoked beside the statute which is challenged and to decide whether the latter squares with the former.''
\end{quote}
\noindent --- Justice Roberts, Opinion of the Court, 297 U.S.\ at 62.

\end{document}
